%% ============================================================
%% PART 2 - Technical Defensive Response
%% ============================================================

\clearpage
\tightsection{Part 2: Technical Defensive Response}
\textbf{Selected risk.} R1, unauthorised privileged access to \texttt{backup01} via suspected compromise of the \texttt{backup} account, reached through opportunistic SSH brute force from \texttt{203.0.113.89} and extended to root through the account's sudo rights.

\textbf{Defensive goal.} The defence targets the full path rather than brute force alone:
\[
\text{SSH brute-force attempts} \rightarrow \text{successful backup login} \rightarrow \text{active session} \rightarrow \text{sudo as root}.
\]
The response is layered: prevention to stop guessing succeeding, containment to limit damage if a credential is still compromised, and detection to catch a success the preventive layer misses.

\textbf{Cyber Kill Chain framing.} The model is represented as a defensive interruption table for the observed SSH path, not as proof that every Cyber Kill Chain stage occurred.

\begin{center}
\scriptsize
\begin{tabularx}{\linewidth}{L{2.2cm}Y L{2.2cm}L{3.35cm}}
\toprule
Kill-chain stage & Representation in this extract & Evidence status & Defensive interruption \\
\midrule
Reconnaissance & Username and service probing through invalid-user and failed-login attempts. & Visible in \texttt{auth.log}. & Detect abnormal probing and rate-limit repeated failures. \\
Delivery / SSH exploitation & Password guessing against exposed SSH reaches a valid \texttt{backup} login. & Visible as failed passwords followed by accepted login. & Key-only SSH, \texttt{PasswordAuthentication no}, fail2ban/rate limiting, and \texttt{PermitRootLogin no}. \\
Installation/C2 & Malware delivery, persistence, and command-and-control are not visible from this extract. & Not proven by \texttt{auth.log}. & Validate with shell history, full journal, audit logs, cron, and file-integrity review. \\
Actions on objectives / privilege action & Accepted \texttt{backup} login is followed by immediate sudo activity as root. & Suspicious privileged activity, not confirmed damage. & Least-privilege sudoers, SELinux confinement, central logging, and sudo/audit logging. \\
\bottomrule
\end{tabularx}
\end{center}

\begin{table}[H]
\centering
\caption{Recommended controls mapped to observed stages.}
\scriptsize
\begin{tabularx}{\linewidth}{L{2.35cm}L{4.65cm}Y}
\toprule
Attack stage & Control & Security effect \\
\midrule
Password guessing & Key-only SSH; \texttt{PasswordAuthentication no}; rate limiting/fail2ban; \texttt{PermitRootLogin no}. & Password guessing no longer grants access by default; automated brute-force noise is reduced and direct root SSH login is closed. \\
Valid-account SSH access & Restrict the \texttt{backup} key by source and forced command; disable interactive shell if not required. & A stolen password is useless; even a stolen key only runs approved backup operations from approved sources. \\
Privilege escalation via sudo & Least-privilege sudoers for \texttt{backup}; no shell, \texttt{su}, or general sudo. & A compromised \texttt{backup} context cannot obtain a general root shell or run arbitrary privileged commands. \\
Post-compromise tampering & SELinux enforcing, with a confined backup domain. & Limits writes to critical system files, persistence locations, and other paths outside the backup role. \\
Evidence gaps after login & Central log forwarding plus sudo/audit logging. & Preserves evidence beyond local \texttt{auth.log} and gives investigators command-level context after authentication. \\
\bottomrule
\end{tabularx}
\end{table}

\textbf{Evidence-backed Defensive Justification.} The controls are tuned to the observed evidence rather than applied as generic SSH hardening.

\begin{center}
\scriptsize
\begin{tabularx}{\linewidth}{L{4.55cm}Y}
\toprule
Part 1 evidence & Defensive use in Part 2 \\
\midrule
\texttt{backup} is the most targeted account, and \texttt{root} is also heavily targeted (Appendix~\ref{app:evidence-risk}, Figure~\ref{fig:app-targeted-users}). & Justifies key-only SSH for \texttt{backup} and \texttt{PermitRootLogin no}, because both the selected asset account and direct root login are prominent targets. \\
\texttt{203.0.113.89} combines high failure volume, accepted \texttt{backup} login, session establishment, and immediate sudo activity (Appendix~\ref{app:raw-excerpt}). & Justifies prioritising R1 over generic brute-force noise: the concern is the complete access-to-privilege path, not the failure count alone. \\
Night traffic shows 3,643 failed logins against 149 accepted logins (Appendix~\ref{app:evidence-risk}, Table~\ref{tab:app-time-window}). & Justifies off-hours weighting in the alert rule, since the 03:08:19 \texttt{backup} login occurs inside the peak attack window. \\
Other mixed-outcome sources exist, but with much lower failure volume and successful legitimate activity. & Justifies threshold and ratio tuning instead of a naive ``failed then accepted'' rule that would produce false positives. \\
The Jul 11 sequence shows accepted login followed by sudo commands, but \texttt{auth.log} cannot show all shell actions. & Justifies least-privilege sudoers, sudo/audit logging, and central log forwarding to improve containment and investigation beyond local authentication records. \\
\bottomrule
\end{tabularx}
\end{center}

\textbf{Prevention: remove the exposed password surface.} The root cause of R1 is password authentication on a privileged, internet-facing account. Disabling \texttt{PasswordAuthentication} and requiring SSH keys removes the entire brute-force success path; fail2ban-style rate-limiting reduces the noise and load from the ongoing automated scanning identified in Part 1. Denying direct root SSH login closes the most-targeted account (\texttt{root}, 1,048 failed attempts) entirely.

\textbf{Containment: least privilege for the \texttt{backup} account.} Because \texttt{backup} is an operational/service account, broad interactive sudo access is unnecessary risk. Restricting its sudoers entry to the exact commands its backup jobs require, and explicitly denying shell, \texttt{su}, and \texttt{sudo}, prevents a compromised credential from being escalated to a general root shell (the T1548.003 step observed at 03:08:19).

\textbf{Containment: SELinux to protect critical files.} Discretionary permissions and sudo rules can still be bypassed if an attacker reaches root. SELinux in enforcing mode adds mandatory access control: the backup service runs in a confined domain that is permitted to touch only its own data paths. Even a compromised \texttt{backup} context is denied write access to sensitive system files such as \texttt{/etc/passwd}, \texttt{/etc/shadow}, \texttt{/etc/sudoers}, cron directories, and SSH host keys, because those files carry protected types the backup domain has no policy allowance for. This limits both privilege escalation (e.g. editing \texttt{/etc/passwd} or \texttt{sudoers}) and anti-forensic tampering (e.g. altering logs or cron persistence), turning a full host compromise into a contained one.

\textbf{Detection: catching a success the preventive layer misses.} If a credential is still compromised, the failure-then-success transition should raise a high-priority alert. A naive ``failed then accepted from the same IP'' rule is too broad; in this extract 13 source IPs match that pattern, but 12 are legitimate users who occasionally mistype before authenticating (13--28 failures, with successes outnumbering failures). Selectivity therefore comes from tuning to volume and ratio, not the raw pattern. The off-hours concentration from Part 1 also motivates a time-aware detection control: the successful \texttt{backup} login occurred at 03:08:19, inside the 00:00--05:59 peak attack window and outside the host's normal login baseline.

Alert when, from a single source IP within a defined window:
\begin{enumerate}
    \item failed SSH attempts exceed a threshold set above the legitimate baseline; legitimate sources here peaked at $\sim$28 failures, so $\geq$50 cleanly separates them (\texttt{203.0.113.89} produced 311);
    \item and the failure-to-success ratio is heavily skewed (legitimate $\approx$ 1:1; \texttt{203.0.113.89} = 311:1);
    \item and the same IP then obtains an accepted login to the privileged \texttt{backup} account (the only account that successfully authenticates in this log; \texttt{root}, \texttt{admin}, and \texttt{deploy} appear only as failed/invalid targets);
    \item weighted higher when the success and follow-on sudo fall outside business hours (the Part 1 attack window was 00:00--05:59; the compromise was at 03:08).
\end{enumerate}

Applied to this extract, this rule flags \texttt{203.0.113.89} and suppresses the 12 legitimate mixed-outcome sources, giving a high-confidence signal that combines volume, ratio, target sensitivity, and timing.

\textbf{Effectiveness, limitations, and residual risk.} Key-only SSH and root-login denial reduce likelihood by removing the password-guessing success path entirely. Least-privilege sudo and SELinux confinement reduce impact by preventing a compromised \texttt{backup} credential from becoming general root control or tampering with critical files. The tuned correlation rule and central logging improve detection and investigation for any success the preventive layer misses.

\textbf{Limitations and trade-offs.}
\begin{itemize}
    \item Detection selectivity depends on tuning. Thresholds set too low re-introduce the 12 legitimate false positives; too high may miss a low-and-slow attacker. The rule mitigates single-source brute force but not a distributed campaign spread across many IPs.
    \item Rate-limiting/lockout can cause denial of service. Aggressive fail2ban thresholds could block legitimate operations or be abused to lock out accounts; thresholds must be tested so scheduled backup jobs are not interrupted.
    \item Key management overhead. Key-only SSH shifts risk to key storage and rotation; a stolen key still authenticates unless source- and command-restricted.
    \item SELinux is not absolute. Policy authoring and maintenance are non-trivial, mislabelling can break services, and an attacker who reaches an unconfined admin role could still relabel or disable enforcement. It raises the cost of compromise rather than eliminating it.
\end{itemize}

\textbf{Residual risk.} The controls improve prevention, containment, detection, and investigation, but they do not retrospectively prove what happened after the Jul 11 login. Confirming or clearing that event still requires the wider evidence review recommended in Part 1 (shell history, full journal, audit logs, cron and file-integrity checks).
